% Generated from locale/tr/content/set-theory/ordinals/wo.tex; do not hand-edit.


\olfileid[tr]{sth}{ordinals}{wo}
\olsection{İyi Sıralamalar}

Temel kavram şöyledir:

\begin{defn}
$<$ bağıntısı, aşağıdaki iki koşulu sağlıyorsa $A$'yı \emph{iyi sıralar}:
\begin{enumerate}
	\item $<$ bağlantılıdır; yani bütün $a,b\in A$ için ya $a<b$ ya $a=b$
	ya da $b<a$ olur;
	\item $A$'nın boş olmayan her alt kümesi bir $<$-minimal
	!!{element} içerir; yani $\emptyset\neq X\subseteq A$ ise
	$(\exists m\in X)(\forall z\in X)z\nless m$ olur.
\end{enumerate}
\end{defn}

Az önce ele aldığımız üç örneğin gerçekten iyi sıralama bağıntıları olduğunu
görmek kolaydır.

\begin{prob}
\Olref[sth][ordinals][idea]{sec}, doğal sayılar üzerinde üç örnek sıralama
sundu. Her birinin iyi sıralama olduğunu denetleyin.
\end{prob}

İyi sıralamaya ilişkin bazı temel fakat son derece önemli gözlemler şunlardır.

\begin{prop}\ollabel{wo:strictorder}
$<$, $A$'yı iyi sıralıyorsa $A$'nın boş olmayan her alt kümesinin biricik
$<$-en küçük elemanı vardır ve $<$ yansımasız, asimetrik ve geçişlidir.
\end{prop}

\begin{proof}
$X$, $A$'nın boş olmayan bir alt kümesiyse bir $<$-minimal !!{element}
$m$ içerir; yani $(\forall z\in X)z\nless m$ olur. $<$ bağlantılı olduğuna
göre $(\forall z\in X)m\leq z$ olur. Böylece $m$, $X$ içinde $<$'ye göre
en küçük !!{element} olur.

Yansımasızlık için $a\in A$ sabitleyin; $\{a\}$ kümesindeki $<$-en küçük
!!{element} $a$ olduğundan $a\nless a$'dır. Geçişlilik için $a<b<c$ ise
$\{a,b,c\}$ kümesinin bir $<$-en küçük !!{element} içermesinden $a<c$
çıkar. Asimetri, yansımasızlık ile geçişlilikten çıkar.
\end{proof}

\begin{prop}\ollabel{propwoinduction}
$<$, $A$'yı iyi sıralıyorsa her $\phi(x)$ formülü için:
\[
	\text{eğer }(\forall a \in A)((\forall b < a)\phi(b) \lif
		\phi(a))\text{ ise }(\forall a \in A)\phi(a).
\]
\end{prop}

\begin{proof}
Karşıt-tersi ispatlayacağız. $\lnot(\forall a\in A)\phi(a)$, yani
$X=\Setabs{x\in A}{\lnot\phi(x)}\neq\emptyset$ olduğunu varsayın. O zaman
$X$ içinde $<$-minimal !!{element}, $a$, vardır. Dolayısıyla
$(\forall b<a)\phi(b)$ olduğu hâlde $\lnot\phi(a)$ olur.
\end{proof}\noindent
Bu son özellik, doğal sayılar üzerindeki güçlü tümevarım ilkesini
hatırlatmalıdır: $(\forall n\in\omega)((\forall m<n)\phi(m)\lif\phi(n))$
ise $(\forall n\in\omega)\phi(n)$ olur. Bu özellik, iyi sıralamayı çok
\emph{sağlam} bir kavrama dönüştürür.\footnote{Hatırlatma: açıkça aksi
belirtilmedikçe bütün formüller parametreler içerebilir.}

